% !TEX root = ../Main.tex
\chapter{双打规则}

\section{发球次序}

\state{发球方的排定}{
    \begin{itemize}
        \item 每盘比赛的\textbf{第 1 局}开始前，发球方的\textbf{两位搭档}协商决定\textbf{由谁先发球}，并在该盘中保持此次序；
        \item 该盘\textbf{第 3 局}时，由先前发球方的\textbf{另一位搭档}发球；
        \item 此后每局由双方队员轮流发球，直到该盘结束。下一盘开始时重新确定发球人。
    \end{itemize}
}

\section{接球次序}

\state{接球方的排定}{
    \begin{itemize}
        \item 每盘比赛\textbf{开始前}，接球方的两位搭档协商决定\textbf{谁接发球手右区的发球（即自己的左区）}、\textbf{谁接发球手左区的发球（即自己的右区）}，整盘均保持此次序；
        \item 当本方发球局结束、轮到本方接球时，\textbf{保持同样的左右站位分工}。
    \end{itemize}
}

\section{回球次序}

\state{发球之后的击球}{
    \begin{itemize}
        \item \textbf{发球}：只能由当局\textbf{指定的发球人}完成；
        \item \textbf{首次接球}：由\textbf{对应站位的接球人}完成（不能由搭档代接）；
        \item \textbf{此后的每一回合}：本方任一搭档\textbf{均可击球}，\textbf{不需要轮流}。
    \end{itemize}
}

\rmkb{
    \textbf{易考点}：
    \begin{itemize}
        \item 发球次序：每盘开始双方队内\textbf{自行约定顺序}，\textbf{在整盘中保持固定}；
        \item 接球次序：\textbf{一盘内左右站位固定}，不能换边接发；
        \item \textbf{首次接发必须由指定人完成}，但\textbf{发球与首次接球之后}，本方任一搭档都可击球；
        \item 每盘开始时顺序可重新商定。
    \end{itemize}
}
